\chapter{Практика}
\label{ch:chap2}

\section{Задание 1: Вычисление ПФ прямоугольного сигнала на основе численного интегрирования}

Сгенерируем одиночный прямоугольный импульс и несущее колебание на частоте 100 Гц

\begin{figure}[H]
    \centering
    \includegraphics{g1.png}
    \caption{Baseband and carrier}
\end{figure}

Перемножим два этих сигнала и получим наш Baseband на частоте 100 Гц. Изобразим его спектр

\begin{figure}[H]
    \centering
    \includegraphics{g2.png}
    \caption{RF signal spectrums}
\end{figure}

В амплитудном спектре видно, что на 100 Гц находится sinc, т.е спектр нашего Baseband сигнала сместился в точку 100 Гц, т.е 
мы промодулировали сигнал. \\

Уменьшим длительность сигнала в несколько раз

\begin{figure}[H]
    \centering
    \includegraphics{g3.png}
    \caption{Spectrum widening}
\end{figure}

Видим, что спектр расширился. Из этого можно сделать вывод, что чем короче импульс, тем шире его спектр.

\section{Задание 2: Спектр последовательности, вычисление при помощи ДПФ}
Заданы следующие параметры

\begin{enumerate}
    \item N=128 - длина битовой последовательности
    \item T=0.01 - длительность символа
    \item ns=64 - кол-во семплов на бит
\end{enumerate}

Определим интервал дискретизации $T_s$. Интервал дискретизации это время, через которое мы будем брать отсчеты. Если на 1 бит приходится
64 семпла, а один бит длится T, то $T_s=\frac{T}{ns}$.\\

Вычислим спектр DFT

\begin{figure}[H]
    \centering
    \includegraphics{g4.png}
    \caption{DFT spectrum, T=0.001s}
\end{figure}

Получили график схожий с sinc.\\

Увеличим T

\begin{figure}[H]
    \centering
    \includegraphics{g5.png}
    \caption{DFT spectrum, T=0.1s}
\end{figure}

Спектр стал более узким. Это еще раз говорит о том, что ширина спектра зависит от длительности импульса. \\

Передадим 20 пакетов и посчитаем средний спектр

\begin{figure}[H]
    \centering
    \includegraphics{g6.png}
    \caption{AVG DFT spectrum, T=0.001s}
\end{figure}

\endinput 